\documentclass[a4paper,11pt,parskip=half-]{scrartcl} 
% scrartcl for Montserrat

\usepackage[a4paper,margin=2cm,marginparwidth=1.75cm]{geometry}

\usepackage[english]{babel}
\usepackage{array}
\usepackage{amsmath}
\usepackage[table]{xcolor}
\usepackage[colorlinks=true, allcolors=blue]{hyperref}
\usepackage{graphicx}

\usepackage[defaultfam,tabular]{montserrat}
\usepackage[scaled]{beramono}
\usepackage[T1]{fontenc}

\usepackage{wrapfig}  % for top-right logo

\usepackage[most]{tcolorbox}
\usepackage{tikz}

\usepackage{caption}
\usepackage{subcaption}

% Custom colours
\definecolor{uwcsblue}{RGB}{58,125,255}
\definecolor{uwcsred}{RGB}{238,79,79}
\definecolor{uwcsyellow}{RGB}{255,199,0}

% For correct/incorrect
\newcommand*\circled[1]{
    \tikz[baseline=(char.base)]{
    \node[shape=circle,inner sep=2pt,fill=uwcsblue] 
    (char) {\textcolor{white}{\textbf{#1}}};
    }
}

% For score-based
\newcommand*\squared[1]{
    \tikz[baseline=(char.base)]{
    \node[shape=rectangle,inner sep=4pt,fill=uwcsred] 
    (char) {\textcolor{white}{\textbf{#1}}};}
}

\usepackage{etoolbox} % for conditional execution
\newtoggle{IsScore}

\newcommand*\problem[4]{
    \addcontentsline{toc}{subsection}{#1}
    {\LARGE 
    \iftoggle{IsScore}
    {
        \squared{#1}
    }
    {
        \circled{#1}
    } 
    \textbf{#2}
    } 
    \hfill
    {\footnotesize Created by: \itshape #3}

    {\itshape #4}
}

\usetikzlibrary{shapes.arrows}

\graphicspath{{../common/}}
\pagestyle{empty}
\setkomafont{section}{\usefont{T1}{fvs}{b}{n}\Large}

\renewcommand{\arraystretch}{1.5}

% Custom logo footer
\usepackage{fancyhdr}

\pagestyle{fancy}
\fancyhf{}
\renewcommand{\headrulewidth}{0pt}
\fancyfoot[EOCF]{\includegraphics[scale=0.05]{logo_wide.png}}
\fancyfoot[R]{\thepage}


\begin{document}

\begin{wrapfigure}{r}{0.2\textwidth}
    \vspace{-10pt} % Move in line with text
    \includegraphics[width=0.9\linewidth]{shield.png}
    \vspace{-100pt} % Prevent text below from moving out of the way
\end{wrapfigure}

{
    \Huge \bfseries UWCS Challenge
}

{
    \Large A friendlier programming competition
}

\section*{Introduction}

Welcome to our internal progcomp!
There will be problems. There will be prizes! And of course, there will be pizza.
This document contains specifications for the 5 problems.

For each problem, you will need to process your solution on various input datasets, then submit the outputs (and your code) for scoring.
This is all handled by our website:
\begin{center}
    \huge \href{https://challenge.uwcs.co.uk}{challenge.uwcs.co.uk}
\end{center}

There are 100 (+25) points available for each problem. 
The base 100 points are divided across each input file. 
For example - if there are 5 input files, each is worth up to 20 points. 
Each input file has multiple test cases, so you can get partial marks.
Fully complete a problem and get a bonus 25 points - a nice incentive for finishing it off!

Later problems may be score-based, so won't have a designated solution.
For these, points are awarded by ranking teams on their score for each input dataset, and scaling that from 0 to 100.
To be ranked, your score must be non-negative.

The first 3 problems will be released simultaneously.
The rest will be staggered. \\ 
Good luck, and have fun!

\vspace{16pt}
\hrule
\vspace{16pt}

\togglefalse{IsScore}
\problem{0}{Reverse}{N/A}
{esreveR}

This sample problem is about reversing strings. \textbf{It does not count towards your score.}

\vspace{8pt}
\hrule

\textbf{Input}

The first line of the input contains an integer $n$ - the number of test cases. \\ 
The next $n$ input lines are strings ($s_1, s_2, \dots, s_n$) of only lowercase characters.

\textbf{Constraints}
\begin{equation*}
    1 \leq n \leq 10^2 \qquad 
    \forall i \in \{1, \dots, n\}, \quad 1 \leq |s_i| \leq 10^2
\end{equation*}

\textbf{Output}

$n$ lines, with line $i$ being the characters of the string $s_i$ in reverse order.

\vspace{8pt}
\hrule

\textbf{Example}

\begin{table}[h]
    \centering
    \begin{tabular}{|p{0.4\linewidth}|p{0.4\linewidth}|}
        \hline
        Input & Output \\
        \hline
        \texttt{\textbf{3}\newline 
        hello\newline 
        goodbye\newline 
        helloagain} & 
        \texttt{olleh\newline
        eybdoog\newline
        niagaolleh} \\
        \hline
    \end{tabular}
\end{table}

\newpage

\togglefalse{IsScore}
\problem{1}{Find Milk}{Keegan R}
{How long does it take for a joke to die?}

The exec were busy preparing for Milk Mondays, the greatest event known to mankind.
But, oh no!
The milk (all four characters of it) got spilt into a pool of letters.

\begin{center}
    \Huge \bfseries \textcolor{uwcsblue}{M}AKESAWH\textcolor{uwcsblue}{I}TE\textcolor{uwcsblue}{L}A\textcolor{uwcsblue}{K}E
\end{center}

Don't cry over spilt milk.
Instead, let's bring it back together: find the shortest substring that contains all of the letters "M", "I", "L", and "K".
Quick, before it goes off!

\vspace{8pt}
\hrule

\textbf{Input}

The first line of the input contains an integer $n$, denoting the number of test cases.

The next $n$ lines each consist of a single string of uppercase letters, $s$.

\textbf{Constraints}
\begin{equation*}
    1 \leq n \leq 10^3 \qquad
    4 \leq |s| \leq 10^4
\end{equation*}

At least one "M", "I", "L", and "K" will be present in $s$.

\textbf{Output}

$n$ lines, each consisting of two space-separated integers: 
both the start and end indices (zero-indexed) of the smallest substring in $s$ containing the characters "M", "I", "L", "K".

\vspace{8pt}
\hrule

\textbf{Example}

\begin{table}[h]
    \centering
    \begin{tabular}{|p{0.4\linewidth}|p{0.4\linewidth}|}
        \hline
        Input & Output \\
        \hline
        \texttt{3\newline
        HIJKLMN\newline
        KLAAAMIIKAL\newline
        MILKMLIKMIKLIMLK\newline}
        & 
        \texttt{1 5\newline 
        5 10\newline
        3 6\newline}
        \\
        \hline
    \end{tabular}
\end{table}

A few notes about these test cases:

\begin{enumerate}
    \item Only the first and last character are excluded. Remember: order doesn't matter.
    \item Although "KLAAAMI" contains MILK, "MIIKAL" is shorter.
    \item Sometimes, there may be multiple valid answers.
\end{enumerate}
\newpage
\togglefalse{IsScore}
\problem{1}{Find Milk}{Keegan R}
{How long does it take for a joke to die?}

The exec were busy preparing for Milk Mondays, the greatest event known to mankind.
But, oh no!
The milk (all four characters of it) got spilt into a pool of letters.

\begin{center}
    \Huge \bfseries \textcolor{uwcsblue}{M}AKESAWH\textcolor{uwcsblue}{I}TE\textcolor{uwcsblue}{L}A\textcolor{uwcsblue}{K}E
\end{center}

Don't cry over spilt milk.
Instead, let's bring it back together: find the shortest substring that contains all of the letters "M", "I", "L", and "K".
Quick, before it goes off!

\vspace{8pt}
\hrule

\textbf{Input}

The first line of the input contains an integer $n$, denoting the number of test cases.

The next $n$ lines each consist of a single string of uppercase letters, $s$.

\textbf{Constraints}
\begin{equation*}
    1 \leq n \leq 10^3 \qquad
    4 \leq |s| \leq 10^4
\end{equation*}

At least one "M", "I", "L", and "K" will be present in $s$.

\textbf{Output}

$n$ lines, each consisting of two space-separated integers: 
both the start and end indices (zero-indexed) of the smallest substring in $s$ containing the characters "M", "I", "L", "K".

\vspace{8pt}
\hrule

\textbf{Example}

\begin{table}[h]
    \centering
    \begin{tabular}{|p{0.4\linewidth}|p{0.4\linewidth}|}
        \hline
        Input & Output \\
        \hline
        \texttt{3\newline
        HIJKLMN\newline
        KLAAAMIIKAL\newline
        MILKMLIKMIKLIMLK\newline}
        & 
        \texttt{1 5\newline 
        5 10\newline
        3 6\newline}
        \\
        \hline
    \end{tabular}
\end{table}

A few notes about these test cases:

\begin{enumerate}
    \item Only the first and last character are excluded. Remember: order doesn't matter.
    \item Although "KLAAAMI" contains MILK, "MIIKAL" is shorter.
    \item Sometimes, there may be multiple valid answers.
\end{enumerate}
\newpage
\togglefalse{IsScore}
\problem{1}{Find Milk}{Keegan R}
{How long does it take for a joke to die?}

The exec were busy preparing for Milk Mondays, the greatest event known to mankind.
But, oh no!
The milk (all four characters of it) got spilt into a pool of letters.

\begin{center}
    \Huge \bfseries \textcolor{uwcsblue}{M}AKESAWH\textcolor{uwcsblue}{I}TE\textcolor{uwcsblue}{L}A\textcolor{uwcsblue}{K}E
\end{center}

Don't cry over spilt milk.
Instead, let's bring it back together: find the shortest substring that contains all of the letters "M", "I", "L", and "K".
Quick, before it goes off!

\vspace{8pt}
\hrule

\textbf{Input}

The first line of the input contains an integer $n$, denoting the number of test cases.

The next $n$ lines each consist of a single string of uppercase letters, $s$.

\textbf{Constraints}
\begin{equation*}
    1 \leq n \leq 10^3 \qquad
    4 \leq |s| \leq 10^4
\end{equation*}

At least one "M", "I", "L", and "K" will be present in $s$.

\textbf{Output}

$n$ lines, each consisting of two space-separated integers: 
both the start and end indices (zero-indexed) of the smallest substring in $s$ containing the characters "M", "I", "L", "K".

\vspace{8pt}
\hrule

\textbf{Example}

\begin{table}[h]
    \centering
    \begin{tabular}{|p{0.4\linewidth}|p{0.4\linewidth}|}
        \hline
        Input & Output \\
        \hline
        \texttt{3\newline
        HIJKLMN\newline
        KLAAAMIIKAL\newline
        MILKMLIKMIKLIMLK\newline}
        & 
        \texttt{1 5\newline 
        5 10\newline
        3 6\newline}
        \\
        \hline
    \end{tabular}
\end{table}

A few notes about these test cases:

\begin{enumerate}
    \item Only the first and last character are excluded. Remember: order doesn't matter.
    \item Although "KLAAAMI" contains MILK, "MIIKAL" is shorter.
    \item Sometimes, there may be multiple valid answers.
\end{enumerate}
\newpage
\togglefalse{IsScore}
\problem{1}{Find Milk}{Keegan R}
{How long does it take for a joke to die?}

The exec were busy preparing for Milk Mondays, the greatest event known to mankind.
But, oh no!
The milk (all four characters of it) got spilt into a pool of letters.

\begin{center}
    \Huge \bfseries \textcolor{uwcsblue}{M}AKESAWH\textcolor{uwcsblue}{I}TE\textcolor{uwcsblue}{L}A\textcolor{uwcsblue}{K}E
\end{center}

Don't cry over spilt milk.
Instead, let's bring it back together: find the shortest substring that contains all of the letters "M", "I", "L", and "K".
Quick, before it goes off!

\vspace{8pt}
\hrule

\textbf{Input}

The first line of the input contains an integer $n$, denoting the number of test cases.

The next $n$ lines each consist of a single string of uppercase letters, $s$.

\textbf{Constraints}
\begin{equation*}
    1 \leq n \leq 10^3 \qquad
    4 \leq |s| \leq 10^4
\end{equation*}

At least one "M", "I", "L", and "K" will be present in $s$.

\textbf{Output}

$n$ lines, each consisting of two space-separated integers: 
both the start and end indices (zero-indexed) of the smallest substring in $s$ containing the characters "M", "I", "L", "K".

\vspace{8pt}
\hrule

\textbf{Example}

\begin{table}[h]
    \centering
    \begin{tabular}{|p{0.4\linewidth}|p{0.4\linewidth}|}
        \hline
        Input & Output \\
        \hline
        \texttt{3\newline
        HIJKLMN\newline
        KLAAAMIIKAL\newline
        MILKMLIKMIKLIMLK\newline}
        & 
        \texttt{1 5\newline 
        5 10\newline
        3 6\newline}
        \\
        \hline
    \end{tabular}
\end{table}

A few notes about these test cases:

\begin{enumerate}
    \item Only the first and last character are excluded. Remember: order doesn't matter.
    \item Although "KLAAAMI" contains MILK, "MIIKAL" is shorter.
    \item Sometimes, there may be multiple valid answers.
\end{enumerate}
\newpage
\togglefalse{IsScore}
\problem{1}{Find Milk}{Keegan R}
{How long does it take for a joke to die?}

The exec were busy preparing for Milk Mondays, the greatest event known to mankind.
But, oh no!
The milk (all four characters of it) got spilt into a pool of letters.

\begin{center}
    \Huge \bfseries \textcolor{uwcsblue}{M}AKESAWH\textcolor{uwcsblue}{I}TE\textcolor{uwcsblue}{L}A\textcolor{uwcsblue}{K}E
\end{center}

Don't cry over spilt milk.
Instead, let's bring it back together: find the shortest substring that contains all of the letters "M", "I", "L", and "K".
Quick, before it goes off!

\vspace{8pt}
\hrule

\textbf{Input}

The first line of the input contains an integer $n$, denoting the number of test cases.

The next $n$ lines each consist of a single string of uppercase letters, $s$.

\textbf{Constraints}
\begin{equation*}
    1 \leq n \leq 10^3 \qquad
    4 \leq |s| \leq 10^4
\end{equation*}

At least one "M", "I", "L", and "K" will be present in $s$.

\textbf{Output}

$n$ lines, each consisting of two space-separated integers: 
both the start and end indices (zero-indexed) of the smallest substring in $s$ containing the characters "M", "I", "L", "K".

\vspace{8pt}
\hrule

\textbf{Example}

\begin{table}[h]
    \centering
    \begin{tabular}{|p{0.4\linewidth}|p{0.4\linewidth}|}
        \hline
        Input & Output \\
        \hline
        \texttt{3\newline
        HIJKLMN\newline
        KLAAAMIIKAL\newline
        MILKMLIKMIKLIMLK\newline}
        & 
        \texttt{1 5\newline 
        5 10\newline
        3 6\newline}
        \\
        \hline
    \end{tabular}
\end{table}

A few notes about these test cases:

\begin{enumerate}
    \item Only the first and last character are excluded. Remember: order doesn't matter.
    \item Although "KLAAAMI" contains MILK, "MIIKAL" is shorter.
    \item Sometimes, there may be multiple valid answers.
\end{enumerate}

\end{document}